
%%%%%%%%%%  scan idx 219  |  folio 212  %%%%%%%%%%


S G A 7


\underline{EXPOSE XVII}


\underline{PINCEAUX DE LEFSCHETZ: THEOREME D'EXISTENCE}


\underline{par N. M. KATZ}


0. \underline{Introduction}

Dans le présent exposé (qui, dans le séminaire oral, se plaçait avant les exposés XIV, XV consacrés à la formule de Picard-Lefschetz), nous introduisons les ``pinceaux de Lefschetz'' de sections hyperplanes d'un schéma projectif et lisse, nous en donnons des conditions d'existence; dans l'exposé suivant, nous en faisons une étude cohomologique élémentaire, basée sur la structure connue de la cohomologie des variétés éclatées. Le fait que certaines hypothèses sur les pinceaux de Lefschetz que nous sommes amenés à introduire sont souvent satisfaites résultera de la formule de Picard-Lefschetz de Exp. XV, de nature moins élémentaire que les résultats du présent exposé et du suivant.

Certains développements du présent exposé et du suivant se trouveront exposés ailleurs dans un cadre plus général (schémas de base généraux, hypothèses de quasi-projectivité au lieu de projectivité, etc.), notamment l'étude pré-cohomologique des pinceaux de sections hyperplanes (et des pinceaux de Lefschetz plus particulièrement) dans EGA V, et l'étude cohomologique des schémas éclatés dans
%%%%%%%%%%  scan idx 220  |  folio 213  %%%%%%%%%%
l'exposé de JOUANOLOU SGA 5 VII; le lecteur trouvera également dans EGA V des démonstrations où la chasse aux diagrammes intrinsèque remplace l'usage des coordonnées, qu'affectionne le rédacteur du présent exposé.

Je tiens à remercier P. DELIGNE et P. GRIFFITHS pour l'aide qu'ils m'ont apportée dans la conception de cet exposé. Dans tout cet exposé, $k$ désigne un corps algébriquement clos, sauf mention expresse du contraire.


1. \underline{Un rappel sur les singularités quadratiques} (cf VI 6).

1.1. \quad Soit $Y$ un $k$-schéma purement de dimension $n$ . Un point fermé $y \in Y$ s'appelle une \underline{singularité quadratique ordinaire} de $Y$ si :

\begin{sloppypar}
a) le complété $\hat{\mathcal{O}}_{Y,y}$ de l'anneau local de $y$ dans $Y$ est isomorphe au quotient $k[[T_1,\ldots,T_n]]/(Q(T))$ des series formelles en $n=1+\dim_k Y$ variables par l'idéal engendré par un élément $Q(T)$ .
\end{sloppypar}

b) $Q$ ne commence qu'en degré deux, et la sous-variété de $\mathbb{P}^{\,n-1}_k$ définie par l'annulation de la partie homogène de degré deux de $Q(T)$ est lisse sur $k$ si $n \geq 2$ , et la partie homogène de $Q(T)$ de degré deux est non-nulle si $n = 1$ .

Si la caractéristique de $k$ est différente de deux, \underline{ou} si $n$ est pair, la condition b) équivaut à la condition
%%%%%%%%%%  scan idx 221  |  folio 214  %%%%%%%%%%
b') $Q$ ne commence qu'en degré deux, et la matrice hessienne de $Q$

\begin{equation*}
\left(\ \frac{\partial^2 Q}{\partial T_i \partial T_j}\,(0)\ \right)_{1 \leq i,j \leq n}
\end{equation*}
est inversible,

Comme on voit en s'appuyant sur le critère jacobienne de lissité.

La  "forme standard "  pour la partie homogène de degré deux $Q_2$ d'une série $Q$ vérifiant b), forme qui peut être obtenue après un changement de variable linéaire, est

\begin{equation}\tag{1.1.1}
\left\{
\begin{array}{l}
Q_2(T) \;=\; \displaystyle\sum_{i=1}^{[n/2]} T_{2i-1}\, T_{2i} \qquad \mathrm{si}\ \ n\ \ \mathrm{est\ pair} \\[4ex]
Q_2(T) \;=\; T_n^2 \;+\; \displaystyle\sum_{i=1}^{\left[\frac{n}{2}\right]} T_{2i-1}\, T_{2i} \quad \mathrm{si}\ \ n\ \ \mathrm{est\ impair,}
\end{array}
\right.
\end{equation}
comme on voit par la théorie des formes quadratique (Bourbaki, Alg.X).

1.2.    Un point fermé $y \in Y$ qui vérifie a), b) s'appelle \underline{singularité quadratique non-dégénérée}. Comme le déterminant de la matrice hessienne est identiquement nul en caractéristique deux si $n$ est impair, les singularités quadratiques non-dégénérées en dimension $n$ n'existent que si, ou bien $n$ est pair, ou bien si $k$ est de caractéristique différente de deux cas dans lesquels ce sont exactement les singularités quadratiques ordinaires.


%%%%%%%%%%  scan idx 222  |  folio 215  %%%%%%%%%%


2. \underline{Les pinceaux de Lefschetz : énoncé des résultats}

2.1.    Désignons par $\mathbf{P}^r$ l'espace projectif à $r$ dimensions sur $k$ , par $\check{\mathbb{P}}^r$ l'espace projectif \underline{dual} des hyperplans dans $\mathbb{P}^r$ , et par $\mathrm{Gr}(1,\ \check{\mathbb{P}}^r)$ la variété des droites dans $\check{\mathbb{P}}^r$ . Un point $D \in \mathrm{Gr}(1,\ \check{\mathbb{P}}^r)$ s'appelle un \underline{pinceau d'hyperplans} dans $\mathbb{P}^r$ . Considérant $D$ comme une droite dans $\check{\mathbb{P}}^r$ , on écrit, par abus de notation, $D = \{H_t\}_{t \in D}$ , les $H_t$ étant les hyperplans de $\mathbb{P}^r$ qui sont les points de la droite $D$ . On appelle \underline{axe de} $D$ la sous-variété linéaire de $\mathbb{P}^r$ de codimension deux qui est l'intersection $H_o \cdot H_\infty$ de deux éléments distincts quelconques du pinceau. Considérant l'axe de $D$ comme un point de $\mathrm{Gr}(r-2,\ \mathbb{P}^r)$ , la variété des sous-variétés linéaires de dimension $r-2$ dans $\mathbb{P}^r$ , la flèche $D \longrightarrow$ l'axe de $D$ n'est autre que l'isomorphisme bien connu d'orthogonalité

\begin{equation*}
\mathrm{Gr}(1,\ \check{\mathbb{P}}^r) \ \ \stackrel{\textstyle\sim}{\longrightarrow} \ \ \mathrm{Gr}(r-2,\ \mathbb{P}^r) \qquad .
\end{equation*}
2.2.    Soit $X$ un $k$-schéma, propre, lisse et irréductible, de dimension $n \geq 1$ , muni d'un plongement $X \hookrightarrow \mathbb{P}^r$ . Un pinceau $D = \{H_t\}_{t \in \mathbb{P}}$ d'hyperplans dans $\mathbb{P}^r$ s'appelle un \underline{pinceau de Lefschetz}   (par rapport au plongement  $X \hookrightarrow \mathbf{P}^r$) si

a) l'axe de $D$ coupe transversalement $X$  (EGA IV 17.13.7);

b) il existe un ouvert dense $U$ de $D$ tel que pour tout $t \in U$ , l'hyperplan $H_t$ coupe transversalement $X$ ;

c) pour $t_o \in D-U$ , $H_{t_o}$ coupe transversalement $X$ sauf en
%%%%%%%%%%  scan idx 223  |  folio 216  %%%%%%%%%%
un point, qui est un point singulier quadratique ordinaire de  $X . H_{t_o}$  .

Nous allons voir plus bas (3.2.1) que les pinceaux de Lefschetz forment un  \underline{ouvert}  de  $\mathrm{Gr}(1, \mathbb{P}^r)$  .

2.3.    Le plongement  $X \hookrightarrow \mathbb{P}^r$  s'appelle un  \underline{plongement de Lefschetz} si les pinceaux de Lefschetz forment un ouvert dense dans $\mathrm{Gr}(1, \check{\mathbb{P}}^r)$ $^{*)}$.

2.4.    Rappellons maintenant la notion de  \underline{multiple d'un plongement donné}  $X \hookrightarrow \mathbb{P}^r$  . C'est le composé du plongement donné  $X \hookrightarrow \mathbb{P}^r$ et du $d$-ième plongement de Segre,  $S_d : \mathbb{P}^r \hookrightarrow \mathbb{P}^{\binom{r+d}{d}-1}$  qui s'écrit en termes des coordonnés homogènes  $X_o,\ldots,X_r$  dans  $\mathbb{P}^r$

\begin{equation*}
(X_o,\ldots,X_r) \xrightarrow{\ \ S_d\ \ } (,\ldots,X^W,\ldots)
\end{equation*}
où  $X^W = X_o^{W_o} \ldots X_r^{W_o}$  parcourt les monômes de degré  $d$  en $X_o,\ldots,X_r$  . Remarquons qu'une "section de  $X$  par une hypersurface de degré  $d$ " $^{**)}$ pour un plongement donné n'est autre qu'une 

$^{*)}$ou plus correctement, dans l'ensemble des points fermés de $\mathrm{Gr}(1,\check{\mathbb{P}}^r)$. Nous nous permettrons par la suite de tels abus de langage.

$^{**)}$En Anglais, on  dit plus commodément: hypersurface section of degree $d$ .


%%%%%%%%%%  scan idx 224  |  folio 217  %%%%%%%%%%

 "section hyperplane" pour  le $d$-ième multiple de ce plongement.

\underline{Théorème 2.5.}    \underline{Soit}  $X$  \underline{irréductible, propre et lisse sur}  $k$  \underline{de dimension} $n \geq 1$  ,  \underline{muni d'un plongement}  $i : X \hookrightarrow P^r$  . \underline{Alors}

2.5.1. \underline{Tout multiple $d$-ième de}  $i$  ,  \underline{avec}  $d \geq 2$  ,  \underline{est un plongement de Lefschetz.}

2.5.2. \underline{Si}  $k$  \underline{est de caractéristique nulle}; \underline{le plongement donné}  $i$ \underline{est un plongement de Lefschetz.}

La démonstration se coupe en plusieurs morceaux, qui suivent.


3. \underline{La variété duale}

3.1.    Soit  $I$  l'Idéal dans  $\mathbb{P}^r$  qui définit  $X$  . On désigne par $N$  le faisceau conormal  $I/I^2$  , faisceau localement libre de rang $r-n$  sur  $X$  , et par  $\mathbb{P}(N)$  le fibré projectif au-dessus de  $X$ formé des droites dans  $\check{N}$  . (Pour  $X = \mathbb{P}^r$  ,  $\mathbb{P}(N)$  est vide.) Le fibré  $\mathbb{P}(N)$  s'identifie à la sous-variété de  $X \times \check{\mathbb{P}}^r$  formée des couples  $(x, H)$  tels que  $H$  soit tangent à  $X$  en  $x$  , par la flèche  $\mathbb{P}(N) \xrightarrow{\ \nu\ } X \times \check{\mathbb{P}}^r$

\begin{equation}\tag{3.1.1}
\mathbb{P}(I/I^2)_x \ni F \xrightarrow{\ \nu\ } (x, \text{l'hyperplan d'équation } \sum \frac{\partial F}{\partial X_i}(x)T_i\ ),
\end{equation}
$F$  désignant une section locale de  $I/I^2$  . On définit

\begin{equation}\tag{3.1.2}
\varphi : \mathbb{P}(N) \longrightarrow \check{\mathbb{P}}^r
\end{equation}
comme étant le composé  $\mathbb{P}(N) \xrightarrow{\ \nu\ } X \times \check{\mathbb{P}}^r \xrightarrow{\ \mathrm{pr}_2\ } \check{\mathbb{P}}^r$ .


%%%%%%%%%%  scan idx 225  |  folio 218  %%%%%%%%%%

3.1.3. On désigne par $\check{X}$ l'image de $\varphi$ ; c'est l'ensemble des hyperplans qui ne coupent pas transversalement $X$ , et s'appelle la \underline{variété duale} de $X$ (plongé dans $\mathbb{P}^r$ par $i$) .

\underline{Proposition} 3.1.4. \underline{La variété duale} $\check{X}$ \underline{est irréductible et propre de dimension} $\leq r-1$ .

\underline{Démonstration.} C'est l'image par $\varphi$ de $\mathbb{P}(N)$ , lui-même irréductible de dimension $r-1$ si $X \neq \mathbb{P}^r$ , et vide si $X = \mathbb{P}^r$ .

\underline{Remarque} 3.1.5. Il est possible de préciser l'interprétation de $\mathbb{P}(N)$ de la façon suivante. Soit $\underline{Y} \subset X \times \check{\mathbb{P}}^r$ le sous-schéma de $X \times \check{\mathbb{P}}^r$ ``d'incidence'' , dont la fibre en tout point $H$ de $\check{\mathbb{P}}^r(k)$ est $X.H$. Il est immédiat que $\underline{Y}$ est lisse de dimension $n+r-1$ , comme fibré projectif de dimension relative $r-1$ sur $X$ , donc $Y \longrightarrow \check{\mathbb{P}}^r$ fait de $\underline{Y}$ une intersection complète relative de dimension relative virtuelle $(n+r-1)-r = n-1$ . Ceci posé, le \underline{sous-schéma jacobien} $\underline{J}_{n-1}(\underline{Y}/\check{\mathbb{P}}^r)$ (VI 5.2 ) \underline{n'est autre que} $\mathbb{P}(N)$ . Grâce à la compatibilité de la formation du sous-schéma jacobien d'une intersection complète relative avec tout changement de base, on en déduit que pour toute droite $D$ dans $\check{\mathbb{P}}^r$ , le sous-schéma jacobien $\underline{J}_{n-1}(\tilde{X}/D)$ , où $\tilde{X} = \underline{Y} \times_{\check{\mathbb{P}}^r} D$ , n'est autre que $\mathbb{P}(N) \times_{\check{\mathbb{P}}^r} D$ .

\underline{Proposition} 3.2. \underline{Le sous-ensemble} $B$ \underline{de} $\check{X}$ \underline{formé des hyperplans} $H$ \underline{tel que} $X \cdot H$ \underline{ait un et un seul point singulier, et que ce point}
%%%%%%%%%%  scan idx 226  |  folio 219  %%%%%%%%%%
\underline{soit quadratique ordinaire, est ouvert} .

C'est un cas particulier de XV 1.3.2, , s'appuyant la théorie des ``modules henséliens'' de R. ELKIK (XV 1.1.4). Le cas où on se restreint aux points quadratiques \underline{non dégénérés} (i.e. où $\mathrm{car}(k) \neq 2$ \underline{ou} $n=\dim X$ est paire)est d'ailleurs évident grâce à 3.3 ci-dessous (dont la démonstration n'utilise pas 3.2).

\underline{Corollaire} 3.2.1. \underline{Sous les hypothèses de} 3.2, \underline{désignons par} $F$ \underline{le fermé} $\check{X}-B$ \underline{de} $\check{X}$ . \underline{Alors les pinceaux de Lefschetz forment un ouvert dans} $\mathrm{Gr}(1, \check{\mathbb{P}}^r)$ , \underline{et les conditions suivantes sont équivalentes:}

3.2.2. \underline{Le plongement donné} $i : X \hookrightarrow \mathbb{P}^r$ \underline{est un plongement de Lefschetz;}

3.2.3. \underline{il existe un pinceau de Lefschetz d'hyperplans (pour le plongement donné).}

3.2.4. \underline{La codimension de} $F$ \underline{dans} $\check{\mathbb{P}}^r$ \underline{est} $\geq 2$ .

\underline{Démonstration.} Une droite $D$ dans $\check{\mathbb{P}}^r$ est un pinceau de Lefschetz si et seulement si elle vérifie:

(3.2.5) L'axe de $D$ coupe transversalement $X$ ,

(3.2.6) $D$ n'est pas contenue dans $X$ ,

(3.2.7) $D$ ne rencontre pas $F$ .


%%%%%%%%%%  scan idx 227  |  folio 220  %%%%%%%%%%

Les conditions (3.2.5) et (3.2.6) définissent des ouverts de $\mathrm{Gr}(1, \check{\mathbb{P}}^{r})$ , et, $F$ étant fermé , il en est de même pour (3.2.7), donc les pinceaux de Lefschetz forment un ouvert dans $\mathrm{Gr}(1, \check{\mathbb{P}}^{r})$ . Parce que $\dim X \leq r-1$ , l'ouvert (3.2.6) est non-vide, et donc dense. L'ouvert (3.2.5) est également non-vide, donc dense. En effet, soit $H$ un hyperplan qui coupe transversalement $X$ ; un tel $H$ existe car $\dim X \leq r-1$ . Parce que $\dim (X \cdot G) \leq r-1$ , il existe un hyperplan $G$ qui coupe transversalement $X \cdot H$ , et par construction la droite $D$ passant par $H$ et $G$ a une axe qui coupe transversalement $X$ . Donc le plongement est de Lefschetz si et seulement si l'ouvert (3.2.7) est non-vide (et donc dense). Or, $F$ étant fermé, (3.2.7) est non-vide si et seulement si $F$ est de codimension $\geq 2$ dans $\check{\mathbb{P}}^{r}$ (cf. EGA V, ou Mumford, Introduction to Algebraic Geometry, p. 88, Cor 3).

\underline{Corollaire} 3.2.8. \underline{Sous les hypothèses de} 3.2., \underline{supposons que} $i : X \hookrightarrow \mathbb{P}^{r}$ \underline{soit un plongement de Lefschetz. Soit} $H_{o}$ \underline{un hyperplan qui coupe transversalement} $X$ , \underline{et désignons par} $\mathrm{Gr}(1, \check{\mathbb{P}}^{r})_{H_{o}}$ \underline{la sous-variété de} $\mathrm{Gr}(1, \check{\mathbb{P}}^{r})$ \underline{formée des droites passant par} $H_{o}$ . \underline{Alors les pinceaux de Lefschetz passant par} $H_{o}$ \underline{forment un ouvert dense dans} $\mathrm{Gr}(1, \check{\mathbb{P}}^{r})_{H_{o}}$ .

\underline{Démonstration.} D'après (3.2.1), les-dits pinceaux forment un ouvert dans $\mathrm{Gr}(1, \check{\mathbb{P}}^{r})_{H_{o}}$ . L'argument qu'on vient de donner pour
%%%%%%%%%%  scan idx 228  |  folio 221  %%%%%%%%%%
établir que les ouverts (3.2.5) et (3.2.6) sont non-vides démontre que leurs intersections avec $\mathrm{Gr}(1, \check{\mathbb{P}}^{r})_{H_{o}}$ sont également non-vides. Il reste à prouver qu'il y a des droites passant par $H_{o}$ qui ne rencontrent pas $F$ . Pour ceci, soit $T$ un hyperplan dans $\check{\mathbb{P}}^{r}$ qui ne contient pas $H_{o}$ . On a un isomorphisme

\begin{center}
\begin{tikzcd}[column sep=huge]
\mathrm{Gr}(1, \check{\mathbb{P}}^{r})_{H_{o}} & T \arrow[l, "\sim"'] \\
\text{la droite passant par } H_{o} \text{ et } \tau & \tau \arrow[l, mapsto]
\end{tikzcd}
\end{center}
% STRUCTURE: 4 nœuds, 2 flèches, disposition 2 lignes x 2 colonnes. Ligne 1 : nœud (1,1) =
%   $\mathrm{Gr}(1, \check{\mathbb{P}}^{r})_{H_{o}}$ ; nœud (1,2) = $T$. Ligne 2 : nœud (2,1) = « la
%   droite passant par $H_{o}$ et $\tau$ » (texte réparti sur deux lignes du typoscript : « la droite
%   passant par » puis « $H_{o}$ et $\tau$ ») ; nœud (2,2) = $\tau$. Flèche 1 : de (1,2) $T$ vers (1,1)
%   $\mathrm{Gr}(1,\check{\mathbb{P}}^{r})_{H_{o}}$, orientée VERS LA GAUCHE, trait long simple,
%   étiquette $\sim$ placée AU-DESSUS du trait (flèche d'isomorphisme). Flèche 2 : de (2,2) $\tau$ vers
%   (2,1), orientée VERS LA GAUCHE, type mapsto (barre verticale du côté de la source, à droite ; pointe
%   à gauche) — soit $\leftarrowtail$ imprimé comme $\longmapsfrom$. Aucune flèche verticale, aucun
%   crochet, aucun pointillé.
donc $\mathrm{Gr}(1, \check{\mathbb{P}}^{r})_{H_{o}}$ ``est'' un espace projectif de dimension $r-1$. Les droites la-dedans passant par un point de $F$ y forment un fermé de dimension $\leq r-2$ , étant l'image du morphisme

\begin{center}
\begin{tikzcd}[column sep=huge]
F \arrow[r] & \mathrm{Gr}(1, \check{\mathbb{P}}^{r})_{H_{o}} \\
f \arrow[r, mapsto] & \text{la droite passant par } H_{o} \text{ et par } f \; .
\end{tikzcd}
\end{center}
% STRUCTURE: 4 nœuds, 2 flèches, disposition 2 lignes x 2 colonnes. Ligne 1 : nœud (1,1) = $F$ ; nœud
%   (1,2) = $\mathrm{Gr}(1, \check{\mathbb{P}}^{r})_{H_{o}}$. Ligne 2 : nœud (2,1) = $f$ ; nœud (2,2) =
%   « la droite passant par $H_{o}$ et par $f$ . » (le point final est imprimé après le $f$). Flèche 1 :
%   de (1,1) $F$ vers (1,2), orientée VERS LA DROITE, trait long simple, SANS étiquette. Flèche 2 : de
%   (2,1) $f$ vers (2,2), orientée VERS LA DROITE, type mapsto (barre verticale à la queue, à gauche ;
%   pointe à droite), SANS étiquette. Aucune flèche verticale, aucun crochet, aucun pointillé, aucune
%   flèche à double pointe.
\underline{Proposition} 3.3. \underline{Le morphisme} $\varphi : \mathbb{P}(N) \longrightarrow \check{\mathbb{P}}^{r}$ \underline{est net (= non ramifié) en un point fermé} $(x_{o}, H)$ \underline{si et seulement si le point} $x_{o}$ \underline{est une singularité quadratique non-dégénérée} (1.2) \underline{de} $X \cdot H$ . \underline{En particulier,} \underline{le sous-ensemble de} $\mathbb{P}(N)$ \underline{formé des points} $(x, H)$ \underline{tel que} $x$ \underline{soit une singularité quadratique non-dégénérée de} $X \cdot H$ \underline{est ouvert.}


%%%%%%%%%%  scan idx 229  |  folio 222  %%%%%%%%%%

\underline{Démonstration}. Nous pouvons nous borner au cas $X \neq \mathbb{P}^r$ . Les propriétés prétendues équivalentes étant indépendantes du choix particulier des coordonnés, on peut les expliciter à l'aide de  coordonnés choisies d'une façon très particulière. Or, $X$ étant une sous-variété lisse de $\mathbb{P}^r$ de dimension $n \geq 1$ , on peut choisir des coordonnés homogènes $X_o,\dots,X_r$ dans $\mathbb{P}^r$ de telle façon que

3.3.1.  le point $x_o \in \mathbb{P}^r$ soit $(1, 0,\dots,0)$ ;

3.3.2.  posant $x_i = X_i/X_o$ pour $i = 1,\dots,r$ , de sorte que le complété $\hat{\underline{O}}_{\mathbb{P}^r,x_o}$ s'identifie à $k[[x_1,\dots,x_r]]$ , l'idéal définissant la trace de $X$ dans $\mathrm{Spec}\ \hat{\underline{O}}_{\mathbb{P}^r,x_o}$ est engendré par des éléments $g_i$ $(i=1,\dots,r-n)$ dans $\underline{O}_{\mathbb{P}^r,x}$ qui, \underline{dans} $\hat{\underline{O}}_{\mathbb{P}^r,\, x_o}$ s'écrivent

\begin{equation*}
g_i = x_{n+i} - h_i(x_1,\dots,x_n) \quad , \quad i = 1,\dots,r-n
\end{equation*}
où les $h_i(x_1,\dots,x_n)$ sont des séries en $x_1,\dots,x_n$ seulement qui se trouvent dans le carré de l'idéal maximal de $k[[x_1,\dots,x_n]]$ .

3.3.3.  l'hyperplan $H$ provient, via (3.1.1) de $g_{n-r}$ .

Il résulte de (3.3.2) que, posant $V = \mathrm{Spec}(\hat{\underline{O}}_{X,x_o})$, il y a une trivialisation

\begin{equation*}
(I/I^2)|V \;\simeq\; \underline{O}_V^{\,r-n}
\end{equation*}
\begin{equation}\tag{3.3.4}
(\lambda_1,\dots,\lambda_{r-n}) \longmapsto \sum_1^{n-r} \lambda_i g_{n+i}
\end{equation}
Il résulte de ceci et de (3.3.3) qu'il existe un voisinage $W$ de $(x, H)$
%%%%%%%%%%  scan idx 230  |  folio 223  %%%%%%%%%%
dans $\mathbb{P}(N)|V$ tel qu'on ait $V \times \mathbb{A}^{r-n-1} \simeq W$ , l'isomorphisme étant donné par

\begin{equation}\tag{3.3.5}
(v,(\lambda_1,\dots,\lambda_{r-n-1})) \in V\times\mathbb{A}^{r-n-1} \longmapsto (v,(g_{r-n}+ \sum_1^{n-r-1} \lambda_i g_i) \in \mathbb{P}(I/I^2)|V
\end{equation}
(où l'on pose $\lambda_i = \wedge_i / \wedge_{n-r}$ ).

En particulier, le complété $\hat{\underline{O}}_{\mathbb{P}(\check{N}),(x,H)}$ de l'anneau local de $\mathbb{P}(N)$ en $(x, H)$ s'identifie à $k[[x_1,\dots,x_n,\lambda_1,\dots,\lambda_{r-n-1}]]$ . Exprimons maintenant $\varphi$ en termes de ces coordonnés, au niveau des complétés:

\begin{equation}\tag{3.3.6}
\varphi(x,\lambda) = \varphi\, (x_1,\dots,x_n,\lambda_1,\dots,\lambda_{r-n-1})
\end{equation}
\begin{equation*}
= \mbox{l'hyperplan d'équation} \quad \sum_{i=0}^{n+r} \varphi_i(x,\lambda)T_i = 0 \quad .
\end{equation*}
Comme l'hyperplan image $\sum \varphi_i\, T_i$ contient le point $(x)$ , point qui s'écrit en coordonnés homogènes

\begin{equation}\tag{3.3.7}
(x) = (1,x_1,\dots,x_n,\, h_1(x),\dots,h_{r-n}(x) \quad ,
\end{equation}
on a

\begin{equation}\tag{3.3.8}
\varphi_o(x,\lambda)= - \sum_{i=1}^{n} x_i\varphi_i(x,\lambda)- \sum_{i=1}^{r-n} h_i(x)\varphi_{n+i}(x,\lambda)
\end{equation}
dans $\hat{\underline{O}}_{\mathbb{P}(N),(x_o,H)}$ et, d'après (3.1.1), pour $i \geq 1$ , posant $D_i = \frac{\partial}{\partial x_i} \in \mathrm{Der}_k(\hat{\underline{O}}_{\mathbb{P}^r,x_o},\, \hat{\underline{O}}_{\mathbb{P}^r,x_o})$ :

\begin{equation}\tag{3.3.9}
\varphi_i(x,\lambda) = D_i(g_{r-n}+ \sum_{i=1}^{r-n-1} \lambda_i\, g_i)\, (x) \quad \mbox{dans} \quad \hat{\underline{O}}_{\mathbb{P}(N),(x_o H)}
\end{equation}

%%%%%%%%%%  scan idx 231  |  folio 224  %%%%%%%%%%

Rappellant  $g_i = x_{n+i} - h_i(x_1,\dots,x_n)$  , on a

\begin{equation}\tag{3.3.10}
\varphi_i(x,\lambda)= -\frac{\partial}{\partial x_i}\left(h_{r-n} + \sum_{i=1}^{r-n-1} \lambda_i h_i\right) \qquad \mbox{pour} \quad 1 \leq i \leq n \;\; ,
\end{equation}
\begin{equation}\tag{3.3.11}
\varphi_i(x,\lambda) = \lambda_i \qquad \mbox{pour} \quad 1 \leq i \leq n \;\; ,
\end{equation}
et

\begin{equation}\tag{3.3.12}
\varphi_r(x,\lambda) = 1
\end{equation}
Désignons par  $m$  l'idéal maximal de  $\hat{\underline{O}}_{\mathbb{P}(N),\,(x_o,H)}$  . Compte tenu de ce que  $h_i \equiv 0 \bmod m^2$  pour  $i = 1,\dots,r-n$  , on trouve

\begin{equation}\tag{3.3.13}
\varphi_o(x,\lambda) \equiv 0 \bmod m^2
\end{equation}
\begin{equation}\tag{3.3.14}
\varphi_i(x,\lambda) \equiv \; -\frac{\partial h_{r-n}}{\partial x_i} \; \bmod m^2 \qquad \mbox{pour} \quad i=1,\dots,n \;\; .
\end{equation}
Mettant ensemble les formules (3.3.11) - (3.3.14), on trouve que la matrice jacobienne de  $\varphi$ en le point  $(x_o,H)$  est 


%%%%%%%%%%  scan idx 232  |  folio 225  %%%%%%%%%%

\begin{equation}\tag{3.3.15}
 \begin{array}{r@{\;}l}
 & \hspace{1.2cm} n \hspace{3.6cm} r-n-1 \\[0.5ex]
\begin{array}{r}
1\,\{ \\[1ex]
n\,\Bigg\{\vphantom{\displaystyle\frac{\partial^2 h_{r-n}}{\partial x_1 \partial x_1}(0)} \\[1.5ex]
\vphantom{\displaystyle\frac{\partial^2 h_{r-n}}{\partial x_1 \partial x_n}(0)} \\[1.5ex]
\\
\\
r-n-1\,\Bigg\{ \\
\\
\\
\\
1\,\{
\end{array}
&
\left(
\begin{array}{ccccccccccc}
0 & & 0 & & 0 & & . & . & . & . & 0 \\[1ex]
-\displaystyle\frac{\partial^2 h_{r-n}}{\partial x_1 \partial x_1}(0) & .\;.\;.\;.\;. & -\displaystyle\frac{\partial^2 h_{r-n}}{\partial x_n \partial x_1}(0) & .\;.\;. & 0 & . & . & . & . & . & 0 \\[2ex]
-\displaystyle\frac{\partial^2 h_{r-n}}{\partial x_1 \partial x_n}(0) & .\;.\;.\;. & -\displaystyle\frac{\partial^2 h_{r-n}}{\partial x_n \partial x_n}(0) & & 0 & & . & . & . & . & 0 \\[2ex]
0 & & 0 & & 1 & 0 & & & & & 0 \\
0 & & . & & 0 & 1 & & 0 & . & . & 0 \\
 & & & & & 0 & & . & & & \\
 & & & & & & & & . & & . \\
 & & & & & & & & & . & 0 \\
0 & & 0 & & 0 & 0 & & . & . & 0 & 1 \\
0 & .\;.\;.\;.\;.\;.\;. & . & & . & . & & . & . & & 0
\end{array}
\right)
\end{array}
\end{equation}
Donc  $\varphi$ est net en  $(x_o,H)$  si et seulement si la matrice hessienne

\begin{equation*}
\left( \displaystyle\frac{\partial^2 h_{r-n}}{\partial x_i \; \partial x_j}(0) \right) \qquad 1 \leq i, \; j \leq n
\end{equation*}
est inversible. D'autre part, (3.3.2) et (3.3.3) nous assurent que le complété de l'anneau local de  $x_o$  en  $X \cdot H$  est

\begin{equation}\tag{3.3.16}
k[[x_1,\dots,x_n]] \,/\, (h_{r-n}(x_1,\dots,x_n)) 
\end{equation}

%%%%%%%%%%  scan idx 233  |  folio 226  %%%%%%%%%%

 et on conclut grâce à la définition (1.2).

\underline{Exemple} 3.4. Supposons que $X$ soit une hypersurface dans $\mathbb{P}^r$ , définie par une équation homogène $F(X_o,\dots,X_r) = 0$ . Alors $\mathbb{P}(N) = X$ , et $\varphi : X \longrightarrow \check{\mathbb{P}}^r$ s'exprime, en coordonnés homogènes, par

\begin{equation}\tag{3.4.1}
\varphi(X_o,\dots,X_r) \;=\; \left( \frac{\partial F}{\partial X_o},\dots, \frac{\partial F}{\partial X_r} \right) \;.
\end{equation}
C'est le \underline{``morphisme de GAUSS''} .

Donnons un exemple où $\varphi$ est partout ramifié (i.e. non net). Pour ceci, soit $p$ la caractéristique de $k$ , $r$ un entier $\geqslant 1$ , et prenons l'hypersurface $X$ d'équation

\begin{equation}\tag{3.4.2}
\sum_{i=0}^{r} X_i^{p^r+1} \;=\; 0 \;,
\end{equation}
pour laquelle on a

\begin{equation}\tag{3.4.3}
\varphi(X_o,\dots,X_r) = (X_o^{p^r},\dots,X_r^{p^r})
\end{equation}
de sorte que $\check{X}$ est reduit à un point si $p = 0$ , et $\check{X} \simeq X$ , avec $\varphi$ le morphisme de Frobenius si $p > 0$ . On observera qu'il résulte de 3.3 (ou de 3.5.0 plus bas) que si $p > 2$ , le plongement envisagé n'est pas un plongement de Lefschetz.

Reprenons le cas général.


%%%%%%%%%%  scan idx 234  |  folio 227  %%%%%%%%%%

\underline{Proposition} 3.5. \underline{Supposons que le morphisme} $\varphi : \mathbb{P}(N) \longrightarrow \check{\mathbb{P}}^r$ \underline{soit} \underline{génériquement net} (i.e. \underline{ne soit pas partout ramifié}). \underline{Alors} $\varphi$ \underline{induit} \underline{un morphisme birationnel de} $\mathbb{P}(N)$ \underline{sur} $\check{X}$ . \underline{Plus précisement, les trois} \underline{ouverts suivants de} $X$ \underline{sont identiques:}

a) \underline{le plus grand ouvert de lissité de} $\check{X}$ ,

b) \underline{le plus grand ouvert} $V$ \underline{de} $\check{X}$ \underline{tel que} $\varphi : \mathbb{P}(N) \longrightarrow \check{X}$ \underline{induise} \underline{un isomorphisme} $\varphi^{-1}(V) \stackrel{\sim}{\longrightarrow} V$ ,

c) \underline{l'ouvert} (cf. 3.2 ) \underline{de} $\check{X}$ \underline{formé des hyperplans} $H$ \underline{tels que} $X \cdot H$ \underline{ait un et un seul point singulier,} \underline{ce dernier étant quadratique non-dégénéré.}

\underline{Corollaire} 3.5.0 \underline{Supposons} car.$k \neq 2$ \underline{ou} $\dim X$ \underline{paire.} \underline{Le plongement} \underline{donné est un plongement de Lefschetz si et seulement si l'une des deux} \underline{conditions est satisfaite:} a) $\varphi$ \underline{est génériquement net,} \underline{ou} b) $\dim \check{X} \leqslant r-2$ .

En effet, la suffisance résulte du critère (3.2.4), car, par 3.5, on a ou $F =$ lieu sing$(\check{X})$ , donc $\dim F < \dim(\check{X}) = r-1$ , ou $\dim(\check{X}) \leqslant r-2$ . Pour la nécessité, on observe que si $D$ est un pinceau de Lefschetz et si $\dim \check{X} = r-1$ , il existe $x \in D \cap \check{X}$ , et par 3.3 et la définition (2.2 c)) $\varphi$ est net en les points de $\varphi^{-1}(x)$ , donc génériquement net, cqfd.

\underline{Démonstration de} 3.5 L'hypothèse implique $\dim \check{X} = \dim \mathbb{P}(N)(=r-1)$ , d'où résulte (3.2) que $\check{X}$ est une hypersurface irréductible dans $\check{\mathbb{P}}^r$ .


%%%%%%%%%%  scan idx 235  |  folio 228  %%%%%%%%%%

Soit $R$ le lieu de ramification de $\varphi$ , qui est par hypothèse un fermé de $\mathbb{P}(N)$ de dimension au plus $r-2$ , de sorte que $\varphi(R)$ est un fermé de $\check{X}$ de dimension au plus $r-2$ . Désignons par $V'$ l'ouvert dense de $\check{X}$

\begin{equation}\tag{3.5.1}
V' = (\check{X} - \varphi(R)) \cap (\check{X} - \text{lieu sing}(\check{X})) ,
\end{equation}
et par $\underline{F}$ son complémentaire dans $\check{X}$

\begin{equation}\tag{3.5.1.1}
\underline{F} = \varphi(R) \cup \text{lieu sing}(\check{X}) .
\end{equation}
Nous allons prouver d'abord que $\varphi$ induit un isomorphisme de $\varphi^{-1}(V')$ avec $V'$ par la méthode naive de construire un inverse $\psi : V' \longrightarrow \varphi^{-1}(V')$. En effet, considérant $\varphi^{-1}(V')$ comme plongé dans $X \times V'$ via (3.1.1) , de sorte que $\varphi$ s'identifie à $\mathrm{pr}_2$ , on voit que la construction de $\psi$ équivaut à la construction d'un morphisme $\psi_1 : V' \longrightarrow \mathbb{P}^r$ tel que le diagramme

\begin{equation}\tag{3.5.2}
\begin{tikzcd}
\mathbb{P}(N) \supset \varphi^{-1}(V') \arrow[r, "\varphi"] \arrow[d] & V' \arrow[d, "\psi_1"] \\
X \arrow[r, hook] & \mathbb{P}^r
\end{tikzcd}
\end{equation}
% STRUCTURE: 4 nodes, 4 arrows. NODES: (1) top-left = $\mathbb{P}(N) \supset \varphi^{-1}(V')$ (a single
%   typed cell: the ambient $\mathbb{P}(N)$, the containment sign $\supset$, then $\varphi^{-1}(V')$);
%   (2) top-right = $V'$; (3) bottom-left = $X$; (4) bottom-right = $\mathbb{P}^r$. ARROWS: (a) top-left
%   --> top-right, horizontal, pointing RIGHT, plain single-headed \longrightarrow, label $\varphi$
%   printed ABOVE the shaft, roughly over its middle; (b) top-left --> bottom-left, vertical, pointing
%   DOWN, plain single-headed arrow, UNLABELLED - its shaft is drawn beneath the '$\mathbb{P}(N)$' part
%   of the top-left cell (not beneath $\varphi^{-1}(V')$), and its head lands just above and to the left
%   of $X$; (c) top-right --> bottom-right, vertical, pointing DOWN, plain single-headed arrow, label
%   $\psi_1$ printed to the RIGHT of the shaft; (d) bottom-left --> bottom-right, horizontal, pointing
%   RIGHT, HOOKED tail at the left end (inclusion / \hookrightarrow), single arrowhead at the right,
%   UNLABELLED. No 2-cells, no dashed arrows, no double-headed arrows, no iso (tilde) marks in this
%   diagram. The tag (3.5.2) is printed in the left margin, level with the middle of the square.
soit commutatif (on pose $\psi(v) = (\psi_1(v), v) \in \mathbb{P}^r \times V'$) .

Pour ceci, remarquons que pour une hypersurface quelconque dans un espace projectif, le morphisme de GAUSS (3.4.1) est ``défini'' en tout point lisse. Appliquons ceci a l'hypersurface $\check{X}$ dans $\check{\mathbb{P}}^r$ ; comme $V'$ est
%%%%%%%%%%  scan idx 236  |  folio 229  %%%%%%%%%%
contenu dans $V$ , l'ouvert de lissité de $\check{X}$ , le morphisme de GAUSS nous donne un morphisme, noté $\psi_1$ :

\begin{equation}\tag{3.5.3}
\psi_1 : V' \longrightarrow \check{\check{\mathbb{P}}}^r \simeq \mathbb{P}^r .
\end{equation}
Il reste à vérifier la commutativité du diagramme (3.5.2). Soit alors $(x_o, H_o)$ un point fermé $\varphi^{-1}(V')$ ; il faut vérifier que $x_o$ , considéré comme étant un hyperplan dans $\check{\mathbb{P}}^r$ , est tangent à $\check{X}$ en $H_o$ , i.e. on va démontrer que pour $(x_o, H_o) \in \varphi^{-1}(V')$ , $H_o$ tangent à $X$ en $x_o \Rightarrow x_o$ tangent à $\check{X}$ en $H_o$ . Comme $\varphi$ est net en $(x_o, H_o)$ , et $\check{X}$ est lisse en $H_o = \varphi(x_o, H_o)$ , $\varphi$ induit un isomorphisme des espaces tangents

\begin{equation}\tag{3.5.4}
T_{\mathbb{P}(N)}((x_o,H_o)) \overset{\varphi}{\overset{\sim}{\longrightarrow}} T_{\check{X}} (H_o) ,
\end{equation}
de sorte que la première partie de la démonstration s'achève par le lemme suivant.

\underline{Lemme} 3.6 (EULER). \underline{Soit} $(x_o, H_o)$ \underline{un point fermé de} $\mathbb{P}(N)$ \underline{, et} \underline{désignons par} $T_L$ \underline{le fibré tangent de l'hyperplan} $L$ \underline{dans} $\check{\mathbb{P}}^r$ \underline{défini par} $x_o$ . \underline{Alors dans} $T_{\check{\mathbb{P}}^r}(H_o)$ \underline{on a}

\begin{equation*}
\varphi (T_{\mathbb{P}(N)}((x_o, H_o))) \subset T_L(H_o) .
\end{equation*}
\underline{Démonstration.} Choisissons des coordonnés homogènes $X_o,\ldots,X_r$ dans $\mathbb{P}^r$ , en termes desquelles $x_o$ devient le point $\underline{A} = (A_o,\ldots,A_r)$ .


%%%%%%%%%%  scan idx 237  |  folio 230  %%%%%%%%%%

Un point de $T_X(x_o)$ est un point de $\mathbb{P}^r$ à valeurs dans $k[\epsilon]/(\epsilon^2)$,

\begin{equation}\tag{3.6.1}
\underline{A} + \epsilon \, \underline{B} \quad ,
\end{equation}
tel que pour toute section locale $H$ de $I$ , on ait

\begin{equation}\tag{3.6.2}
\sum_{i=0}^{r} \; B_i \, \frac{\partial H}{\partial X_i}(\underline{A}) \;=\; 0 \quad .
\end{equation}
Choisissons de plus une section locale $H$ de $I$ qui donne lieu à $H_o$ par (3.1.1). Par la trivialité locale de $\mathbb{P}(N)$ au-dessus de $X$ , un point de $T_{\mathbb{P}(N)}((x_o,H_o))$ se représente par un point

\begin{equation}\tag{3.6.3}
(\underline{A} + \epsilon \, \underline{B} \, , \; H + \epsilon \, F) \quad ,
\end{equation}
où $F$ est une section locale de $I$ , et où $\underline{A}$ et $\underline{B}$ sont comme ci-dessus. L'image par $\varphi$ de (3.6.3) est l'hyperplan à coefficients dans $k[e]/(\epsilon^2)$ d'équation (c.f. 3.1.1)

\begin{equation}\tag{3.6.4}
\sum_{i=0}^{r} \; \frac{\partial}{\partial X_i} \, (H+\epsilon F) \; (\underline{A}+\epsilon\underline{B})T_i \quad ,
\end{equation}
de sorte que le lemme équivaut à la formule

\begin{equation}\tag{3.6.5}
\sum_{i=0}^{r} \; A_i \, \frac{\partial}{\partial X_i} \, (H+\epsilon F) \; (\underline{A}+\epsilon\underline{B}) = 0 \quad .
\end{equation}
Compte tenu de ce que $\epsilon^2 = 0$ , et la formule de Taylor, on a

\begin{equation}\tag{3.6.6}
\frac{\partial}{\partial X_i}(H+\epsilon F)(\underline{A}+\epsilon\underline{B}) = \frac{\partial H}{\partial X_i}(\underline{A})+\epsilon \, \frac{\partial F}{\partial X_i}(\underline{A})+\epsilon \sum_{j=0}^{r} B_j \, \frac{\partial^2 H}{\partial X_j \partial X_i} \, (\underline{A}) \quad , \quad 
\end{equation}

%%%%%%%%%%  scan idx 238  |  folio 231  %%%%%%%%%%

 donc il suffit de vérifier les trois formules suivantes:

\begin{equation}\tag{3.6.7}
\sum_{i=0}^{r} \; A_i \, \frac{\partial H}{\partial X_i}(\underline{A}) \;=\; 0 \quad ,
\end{equation}
\begin{equation}\tag{3.6.8}
\sum_{i=0}^{r} \; A_i \, \frac{\partial F}{\partial X_i}(\underline{A}) \;=\; 0 \quad ,
\end{equation}
\begin{equation}\tag{3.6.9}
\sum_{i,j=0}^{r} \; A_i B_j \, \frac{\partial^2 H}{\partial X_i \partial X_j} \; (\underline{A}) \;=\; 0 \quad .
\end{equation}
Or, $H$ et $F$ étant homogènes de degré zéro, la formule d'Euler donne

\begin{equation}\tag{3.6.10}
\sum X_i \, \frac{\partial H}{\partial X_i} = 0 \quad , \quad \sum \; X_i \, \frac{\partial F}{\partial X_i} \;=\; 0
\end{equation}
et, $\frac{\partial H}{\partial X_j}$ étant homogène de degré $-1$ , cette formule donne

\begin{equation*}
\sum_i \; X_i \, \frac{\partial^2 H}{\partial X_i \partial X_j} \;=\; - \, \frac{\partial H}{\partial X_j} \quad ,
\end{equation*}
de sorte que (3.6.9) se récrit

\begin{equation}\tag{3.6.11}
- \sum_j B_j \, \frac{\partial H}{\partial X_j}(\underline{A}) = 0 \quad ,
\end{equation}
qui n'est autre que (3.6.2).

Ceci démontre que $\varphi$ induit un isomorphisme

\begin{equation}\tag{3.6.12}
\varphi^{-1}(V') \; \stackrel{\sim}{\longrightarrow} \; V' \quad 
\end{equation}

%%%%%%%%%%  scan idx 239  |  folio 232  %%%%%%%%%%

 dont l'inverse est $\psi$ . Remarquons maintenant que le morphisme

\begin{equation}\tag{3.6.13}
\left\{ \begin{array}{l} \psi \ : \ V' \longrightarrow \ \mathbb{P}^r \times V' \subset \mathbb{P}^r \times \check{X} \\[2ex] \psi(v) \ = \ (\psi_1(v),v) \end{array} \right.
\end{equation}
s'étend en un morphisme

\begin{equation}\tag{3.6.14}
\left\{ \begin{array}{l} \widetilde{\psi} \ : \ V \longrightarrow \ \mathbb{P}^r \times V \\[2ex] \widetilde{\psi}(v) \ = \ (\widetilde{\psi}_1(v),v) \ \ , \end{array} \right.
\end{equation}
où l'on désigne par

\begin{equation}\tag{3.6.15}
\widetilde{\psi}_1 \ : \ V \longrightarrow \ \mathbb{P}^r
\end{equation}
le morphisme de GAUSS, vu comme défini sur $V$ tout entier.

On vient de démontrer l'inclusion

\begin{equation}\tag{3.6.16}
\psi(V') \ \subset \ \mathbb{P}(N)
\end{equation}
ce qui implique, $V'$ étant dense dans $V$ , et $\mathbb{P}(N)$ étant fermé dans $\mathbb{P}^r \times \check{X}$ , l'inclusion

\begin{equation}\tag{3.6.17}
\widetilde{\psi}(V) \subset \mathbb{P}(N)
\end{equation}
Il est évident qu'on a même

\begin{equation}\tag{3.6.18}
\widetilde{\psi}(V) \ \subset \ \varphi^{-1}(V) \ .
\end{equation}
Encore par "continuité" , les relations 


%%%%%%%%%%  scan idx 240  |  folio 233  %%%%%%%%%%

\begin{equation*}
 \left\{ \begin{array}{l} \psi \ \circ \ \varphi \ = \ \mathrm{id}_{\varphi^{-1}(V')} \\[3ex] \varphi \ \circ \ \psi \ = \ \mathrm{id}_{V'} \end{array} \right.
\end{equation*}
impliquent les relations

\begin{equation*}
\left\{ \begin{array}{l} \widetilde{\psi} \ \circ \ \varphi = \mathrm{id}_{\varphi^{-1}(V)} \\[3ex] \varphi \ \circ \ \widetilde{\psi} = \mathrm{id}_{V} \qquad , \end{array} \right.
\end{equation*}
ce qui démontre que les ouverts 3.5 a) et 3.5 b) sont identiques.

Démontrons que les ouverts 3.5 a) et 3.5 c) sont identiques. On remarque d'abord que, d'après 3.3, l'ouvert 3.5 c) est le sous-ensemble $\underline{U}$ de $\check{X}$ défini par

\begin{equation}\tag{3.6.19}
\underline{U} = \left\{\, x \in \check{X} \ \middle|\ \begin{array}{l} \varphi^{-1}(x) \ \mbox{contient exactement un point,} \\ \mbox{et en ce point } \varphi \mbox{ est net.} \end{array} \right\} \ .
\end{equation}
Or $\varphi^{-1}(\underline{U})$ est normal (étant lisse) et le morphisme induit

\begin{equation*}
\varphi \ : \ \varphi^{-1}(\underline{U}) \longrightarrow \underline{U}
\end{equation*}
est birationnel (on l'a déja démontré), radiciel (par définition de $\underline{U}$) et propre. Donc

\begin{equation*}
\varphi \ : \ \varphi^{-1}(\underline{U}) \longrightarrow \underline{U}
\end{equation*}
est le \underline{normalisé} de $\underline{U}$ . En particulier, $\underline{U}$ est \underline{géométriquement unibranche} . Comme $\varphi : \varphi^{-1}(\underline{U}) \longrightarrow \underline{U}$ est \underline{net} , il s'en déduit (EGA IV 18.10.1) qu'il est \underline{étale} . Par EGA IV 18.10.18, on conclut alors que c'est une immersion ouverte, donc un isomorphisme, donc que
%%%%%%%%%%  scan idx 241  |  folio 234  %%%%%%%%%%
$\underline{U}$ est lisse, i.e. qu'on a l'inclusion 3.5 c) $\subset$ 3.5 a). D'autre part, (3.6.19) rend évidente l'inclusion 3.5 b) $\subset$ 3.5 c), et on conclut, compte tenu de ce que 3.5 a)$=$ 3.5 b).

Nous pouvons maintenant démontrer:

3.7.  Le théorème 2.5 est vrai dans le cas $\mathrm{car}.k \neq 2$ , ainsi que dans le cas $X$ de dimension paire.

\underline{Démonstration}.  Compte tenu de 3.5.0 et de 3.3, il suffit de vérifier :

(3.7.1) Pour tout point fermé $x_o \in X$, et tout entier $d \geqslant 2$, il existe une hypersurface $H$ de degré $d$ tangente à $X$ en $x_o$ , telle que $x_o$ soit une singularité quadratique ordinaire de $X.H$.

Or, prenant convenablement des coordonnés homogènes $X_o,\ldots,X_r$ , on peut supposer que $x_o$ est le point $(1,0,\ldots,0)$ , et que les fonctions $X_i/X_o$ , $i = 1,\ldots,n$ , donnent des coordonnés locales sur $X$ dans un voisinage de $x_o$ . On prend pour $H$ l'hypersurface d'équation

\begin{equation}\tag{3.7.1.1}
H = X_o^{d-2} \sum_{i=1}^{n} X_i^2 \qquad \text{si} \quad \mathrm{car}.(k) \neq 2
\end{equation}
\begin{equation}\tag{3.7.1.2}
H = X_o^{d-2} \sum_{i=1}^{n/2} X_i \, X_{n/2+i} \qquad \text{si} \quad n \ \text{est pair}
\end{equation}
\begin{equation}\tag{3.7.1.3}
H = X_o^{d-2} \, (X_n^2 + \sum_{i=1}^{[n/2]} X_i \, X_{[n/2]+i}) \ \text{si} \quad n \ \text{est impair} \ .
\end{equation}

%%%%%%%%%%  scan idx 242  |  folio 235  %%%%%%%%%%

Pour l'énoncé 2.5.2, remarquons qu'ou bien $\varphi$ est génériquement net, et 3.5.0 s'applique, ou bien $\varphi$ est partout ramifié, ce qui implique $\dim X \leqslant r-2$ , $k$ étant de caractéristique nulle. Dans ce dernier cas, le critère (3.2.4) est trivialement vérifié.


4.  \underline{Le cas général}

Nous allons esquisser une autre méthode de démonstration de 2.5 valable en toute caractéristique, en admettant l'énoncé suivant, voisin de 3.1, et cas particulier de XV 1.1.4.

(4.1) \underline{Le sous-ensemble de} $\mathbb{P}(N)$ \underline{formé des points} $(x, H)$ \underline{tels que} $x$ \underline{soit une singularité quadratique ordinaire de} $X \cdot H$ \underline{est ouvert.}

Remarquons que 4.1 résulte de 3.3, sauf si $\mathrm{car}(k) = 2$ \underline{et} $\dim X$ impaire.

4.1.1. Désignons par $R_1$ le fermé de $\mathbb{P}(N)$ complément de l'ouvert 4.1, et par $F_1$ son image $\varphi(R_1) \subset \check{X} \subset \check{\mathbb{P}}^r$ .

\underline{Lemme} 4.1.2. \underline{Si l'ouvert} 4.1 \underline{est non-vide,} \underline{alors} $\operatorname{codim}_{\check{\mathbb{P}}^r}(F_1) \geqslant 2$

\underline{Démonstration}. Par hypothèse, $\dim(R_1) < \dim(\mathbb{P}(N)) = r-1$ , et $F_1 = \varphi(R_1)$ .

4.1.3. Désignons par $F_2$ la partie de $\check{X}$ formée des hyperplans qui sont tangents à $X$ en au moins deux points. On remarque que $F_2$ n'est
%%%%%%%%%%  scan idx 243  |  folio 236  %%%%%%%%%%
pas fermé, mais qu'on a

l'adhérence de $F_2 \subset F_1 \cup F_2$ ,

ceci à cause du fait que $F_1 \cup F_2 = F$ (cf. 3.2) est fermé. (L'inclusion 4.1.4. ``dit'' qu'une \underline{confluence} de deux singularités n'est pas quadratique ordinaire.) En tout cas la construction de 4.2.4. montrera que $F_2$ est constructible. Pour démontrer 2.5.1 nous allons appliquer le critère 3.2.4. Pour ceci, il faudra savoir que $\mathrm{codim}_{\check{\mathbb{P}}^r}(F_1) \geq 2$ et que $\mathrm{codim}_{\check{\mathbb{P}}^r}(F_2) \geq 2$. D'autre part, 3.8.1 et 4.1.2 nous assurent que, remplaçant le plongement donné par son $d$-ième multiple, $d \geq 2$ , on aura $\mathrm{codim}_{\check{\mathbb{P}}^r}(F_2) \geq 2$ . On conclut que 2.5.1 sera prouvé, une fois connue la

\underline{Proposition} 4.2. \underline{Soit} $X$ \underline{une sous-variété lisse et irréductible de} $\mathbb{P}^r$ , \underline{de dimension} $n \geq 1$ . \underline{Pour tout entier} $d \geq 2$ , \underline{la partie} $F_2$ \underline{de} $\check{\mathbb{P}}^{{n+d \choose d}-1} = \check{\mathbb{P}}^N$ \underline{formée des hypersurfaces de degré} $d$ \underline{dans} $\mathbb{P}^r$ \underline{qui sont tangentes à} $X$ \underline{en au moins deux points est de codimension} $\geq 2$ \underline{dans} $\check{\mathbb{P}}^N$ .

4.2.1. Commencons par le cas spécial $d = 2$ et $X$ une sous-variété linéaire de $\mathbb{P}^r$ . Or, il suffira de trouver une droite dans $\check{\mathbb{P}}^N$ qui ne rencontre pas $F$ (cf. 3.2.4). Choisissons des coordonnés homogènes $X_o,\ldots,X_r$ dans $\mathbb{P}^r$ de telle facon que $X$ soit défini par
%%%%%%%%%%  scan idx 244  |  folio 237  %%%%%%%%%%
$X_{n+1} = \ldots = X_r = 0$ , et des coordonnés homogènes $(\lambda, \mu)$ dans $\mathbb{P}^1$ .

Pour $n$ pair, $n = 2k$ , on prend le pinceau de quadriques

\begin{equation}\tag{4.2.2}
H_{(\lambda,\mu)} = \lambda\ X_o^2 + (\lambda X_k + \mu X_o) X_{2k} + \sum_{i=1}^{k-1}\ (\lambda X_i + \mu X_{i+1}) X_{k+i} \quad ;
\end{equation}
pour $n$ impair, $n = 2k+1$ , on prend le pinceau

\begin{equation}\tag{4.2.3}
H_{(\lambda,\mu)} = (\lambda X_1 + \mu X_o) X_o + \sum_{i=1}^{k}\ (\lambda X_{2i} + \mu X_i) X_{2i+1} \quad .
\end{equation}
Dans les deux cas, on constate que $X \cdot H_{(\lambda,\mu)}$ est lisse pour $(\lambda,\mu) \neq (0,1)$ , et que $X \cdot H_{(0,1)}$ n'a qu'un point singulier, qui est même quadratique ordinaire, et donc que le pinceau ne coupe pas $F$ .

4.2.4. Désignons par $\mathrm{Btg}^d(X)$ la sous-variété de $\check{\mathbb{P}}^N \times (X \times X - \Delta(X \times X))$ formée des triples $(H,\ x_1,\ x_2)$ tels que $H$ est une hypersurface de degré $d$ dans $\mathbb{P}^r$ qui est tangente à $X$ en $x_1$ et en $x_2$ , et $x_1 \neq x_2$ . L'image de $\mathrm{Btg}^d(X)$ par $\mathrm{pr}_1$ est l'ensemble $F_2$ de (4.2), de sorte qu'il suffira de démontrer

$\dim \mathrm{Btg}^d(X) \leq N-2$ si $d \geq 3$ , ou si $d = 2$ et $X$ n'est pas linéaire.

Pour ceci, désignons par $\mathrm{Lin}(X \times X)$ la sous-variété de $X \times X - \Delta(X \times X)$

$\mathrm{Lin}(X \times X) = \{(x_1,x_2) \in X \times X \mid x_1 \neq x_2$ , et la droite engendrée par $x_1,x_2$ est tangente à $X$ en $x_1$ $\}$ \\ $\{(x_1,x_2) \in X \times X \mid x_1 \neq x_2$ et $x_2 \in X \cap \mathrm{Tg}\ (x_1)\}$ .


%%%%%%%%%%  scan idx 245  |  folio 238  %%%%%%%%%%

On a évidemment :

4.2.7.\quad \underline{Si} $X$ \underline{n'est pas une sous-variété linéaire de} $\mathbb{P}^r$ , \underline{alors}

\begin{equation*}
\dim \mathrm{Lin}(X \times X) < 2\dim X \quad .
\end{equation*}
Donc pour (4.2.5) il suffira de démontrer que les fibres du morphisme

\begin{equation}\tag{4.2.8}
\mathrm{Btg}^d(X) \stackrel{\mathrm{pr}_{2,3}}{\longrightarrow} X \times X - \Delta(X \times X) \quad ,
\end{equation}
considérés comme des sous-variétés de $\check{\mathbb{P}}^N$ , sont

\begin{equation}\tag{4.2.9}
\left\{
\begin{array}{l}
\mbox{de codimension} \geq 2\dim(X)+1 \quad \mbox{au-dessus de} \quad \mathrm{Lin}(X{\times}X) \quad , \\[2ex]
\mbox{de codimension} \geq 2\dim(X)+2 \quad \mbox{au-dessus de} \quad X{\times}X-\Delta(X{\times}X)-\mathrm{Lin}(X{\times}X)
\end{array}
\right\} \quad .
\end{equation}
L'énoncé (4.2.9) est une conséquence du lemme suivant, en prenant pour $L_i$ l'espace tangent à $X$ en $x_i$ , $i=1,2$ .

\underline{Proposition} 4.3.\quad \underline{Soient} $x_1$ \underline{et} $x_2$ \underline{deux points distincts fermés de} $\mathbb{P}^r$ , \underline{et} $L_1$ \underline{et} $L_2$ \underline{deux sous-espaces linéaires de} $\mathbb{P}^r$ \underline{tels que} $x_1 \in L_1$ \underline{et} $x_2 \in L_2$ . \underline{Dans l'espace} $\check{\mathbb{P}}^N = \check{\mathbb{P}}^{\binom{r+d}{d}-d}$ \underline{des hypersurfaces de degré} $d$ , \underline{le sous-espace linéaire formé des hypersurfaces tangentes à} $L_1$ \underline{en} $x_1$ \underline{et à} $L_2$ \underline{en} $x_2$ \underline{est de codimension}

\begin{equation*}
\left\{
\begin{array}{l}
= \dim L_1+\dim L_2+2 \quad \mbox{si } d\geq 3, \mbox{ ou si } d=2 \mbox{ et } L_1 \cap L_2 \mbox{ ne contient pas} \\[1ex]
\qquad\qquad\qquad\qquad \mbox{la droite } (x_1,x_2) \mbox{ contenant } x_1,x_2 \\[2ex]
= \dim L_1+\dim L_2+1 \quad \mbox{si } d=2 \mbox{ et } L_1 \cap L_2 \mbox{ contient la droite } (x_1,x_2).
\end{array}
\right.
\end{equation*}

%%%%%%%%%%  scan idx 246  |  folio 239  %%%%%%%%%%

\underline{Démonstration.}\quad Pour $d \geq 3$ , il faut prouver que $\dim L_1 + \dim L_2+2$ conditions linéaires sont indépendantes , de sorte qu'il suffit de traiter le cas $L_1 = L_2 = \mathbb{P}^r$ , i.e. de trouver ``combien'' d'hypersurfaces sont singulières en $x_1$ et en $x_2$ . Or, en termes de coordonnés homogènes choisies de telle façon que $x_1 = (1,0,...0)$ et $x_2 = (0,1,0...,0)$ , l'hypersurface d'équation $\Sigma A_W X^W$ est singulière en $(1,0,..,0)$ si et seulement si les coéfficients des monômes $X_o^{d-1}X_i$ , $i=0,\ldots,r$, sont tous nuls, et elle est singulière en $(0,1,0,..0)$ si et seulement si les coefficients des monômes $X_1^{d-1}X_i$ , $i=0,\ldots,r$ , sont tous nuls. Pour $d \geq 3$ , les monômes $X_o^{d-1}X_i$, $i=0,\ldots,r$ et $X_1^{d-1}X_i$, $i=0,\ldots,r$ sont tous distincts, ce qui donne l'énoncé voulu pour $d \geq 3$ .

Pour $d = 2$ , le même calcul donne l'inégalité

\begin{equation*}
\mathrm{codim} \geq \dim L_1 + \dim L_2 + 1 \quad ,
\end{equation*}
(car le monôme $X_o \, X_1$ est compté deux fois).

Vérifions que si $L_1 \cap L_2$ ne contient pas la droite $(x_1,x_2)$, alors les conditions de tangentialité sont indépendantes. On se ramène au cas $L_1 = \mathbb{P}^r$ , $L_2$ un hyperplan dans $\mathbb{P}^r$ qui ne contient pas la droite $(x_1,x_2)$ . On choisit des coordonnés homogènes $X_o,\ldots,X_r$ de telle façon que $x_1 = (1,0,\ldots,0)$ , $x_1 = (0,1,0...0)$ et $L_2$ soit l'hyperplan $X_o = 0$ . Or, une quadrique $\sum\limits_{i \leq j} A_{ij} \, X_i \, X_j$ est singulière en $(1,0,..,0)$ si et seulement si $A_{oj} = 0$ pour $j = 0,\ldots,r$ , et tangente à
%%%%%%%%%%  scan idx 247  |  folio 240  %%%%%%%%%%
$X_o = 0$ en $(0,1,\ldots,0)$ si et seulement si $A_{ij} = 0$ pour $j = 1,\ldots,r$ , ce qui fait ensemble $2r+1 = \dim L_1 + \dim L_2 + 2$ conditions indépendantes.

Pour le cas où $L_1 \cap L_2$ contient la droite $(x_1,x_2)$, on se ramène au cas $L_1 = L_2 =$ la droite $(x_1,x_2)$ . En termes de coordonnés homogènes dans lesquelles $x_1 = (1,0,..,0)$ , $x_2 = (0,1,0,...0)$ , la quadrique $\sum_{i \leq j} A_{ij} X_i X_j$ passe par $(1,0,..0)$ et y est tangente à $L_1$ si et seulement si $A_{oo} = A_{o1} = 0$ ; elle passe par $(0,1,0...0)$ et y est tangente à $L_2$ si et seulement si $A_{11} = A_{o,1} = 0$ , ce qui ne fait ensemble que $3 = \dim L_1 + \dim L_2 + 1$ conditions.


5.\quad \underline{Le degré de la variété duale par voie "élémentaire"}

(Cf.XVIII 3.2. pour une autre méthode dans le cas où $\dim X$ est pair \underline{ou} $\mathrm{car}(k) \neq 2$ .)

Nous employons toujours les notations de 2.2, donc $X \longrightarrow \mathbb{P}^r$ est une sous-variété lisse irréductible de dimension $n$ .

5.1. Commençons par considérer le morphisme (3.1.1)

\begin{equation}\tag{5.1.0}
\nu : \mathbb{P}(N) \longrightarrow X \times \check{\mathbb{P}}^r
\end{equation}
d'un point de vue plus "fonctoriel" . (Je tiens à remercier S. KLEIMAN pour m'avoir signalé la formulation qui suit.) On a une suite exacte
%%%%%%%%%%  scan idx 248  |  folio 241  %%%%%%%%%%
de faisceaux cohérents sur $X$

\begin{equation}\tag{5.1.1}
0 \longrightarrow \check{N} \longrightarrow \Omega^1_{\mathbb{P}/k} \,|\, X \longrightarrow \Omega^1_X \longrightarrow 0 \quad .
\end{equation}
D'autre part, "écrivant les formes différentielles sur $\mathbb{P}^r$ en coordonnés homogènes" , on trouve une suite exacte de faisceaux sur $\mathbb{P}^r$

\begin{equation}\tag{5.1.2}
0 \longrightarrow \Omega^1_{\mathbb{P}^r/k} \longrightarrow (\underline{O}_{\mathbb{P}}(-1))^{r+1} \longrightarrow \underline{O}_{\mathbb{P}} \longrightarrow 0
\end{equation}
qui, restreinte à $X$ , donne alors une suite exacte

\begin{equation}\tag{5.1.3}
0 \longrightarrow \Omega^1_{\mathbb{P}^r/k} \,|\, X \longrightarrow (\underline{O}_X(-1))^{r+1} \longrightarrow \underline{O}_{\mathbb{P}} \longrightarrow 0 \quad .
\end{equation}
Mettant ensemble 5.1.1 et 5.1.3, on trouve une inclusion

\begin{equation}\tag{5.1.4}
\check{N} \lhook\joinrel\longrightarrow (\underline{O}_X(-1))^{r+1} \quad ,
\end{equation}
qui donne, par transposition, une surjection,

\begin{equation}\tag{5.1.5}
\underline{O}_X(1)^{r+1} \longrightarrow N
\end{equation}
d'où une immersion fermée

\begin{equation}\tag{5.1.6}
\alpha : \mathbb{P}(N) \longrightarrow \mathbb{P}(\underline{O}_X(1)^{r+1}) \;\simeq\; X \times \check{\mathbb{P}}^r
\end{equation}
qui n'est autre que (5.1.0).


%%%%%%%%%%  scan idx 249  |  folio 242  %%%%%%%%%%


5.2.  \underline{L'idée du calcul}.

Nous supposons à partir de maintenant que $\check{X} = \varphi(\mathbb{P}(N))$ est de dimension $r-1$ , donc que c'est une hypersurface dans $\check{\mathbb{P}}^r$ . Le calcul de son degré se fait naturellement en introduisant les classes de Chern [1] .

Pour toute variété projective et lisse $T$ , nous notons par $\mathrm{CH}(T)$ son anneau de Chow , qui est un anneau gradué (par la codimension d'un cycle). Etant donné un élément $z \in \mathrm{CH}(T)$ , nous notons par

\begin{equation*}
\int_T z
\end{equation*}
le degré de sa composante de degré zéro (qui est un $0$-cycle). La classe dans $\mathrm{CH}^1(T)$ d'un faisceau inversible $\underline{L}$ sur $T$ sera noté également $\underline{L}$ quand aucun confusion n'est à craindre (par exemple, quand elle se trouve sous le signe d'intégration $\int_T$ ) .

Ceci posé, notant par $L$ la classe dans $\mathrm{CH}(\check{\mathbb{P}}^r)$ d'une droite, on a

\begin{equation}\tag{5.2.1}
\deg(\check{X}) = \int_{\check{\mathbb{P}}^r} L\cdot\check{X} \quad .
\end{equation}
Or, $\check{X}$ est ``ensemblistement'' l'image de $\mathbb{P}(N)$ par $\varphi$ (3.1.2). Mais, du point de vue de l'homomorphisme ``image directe''

\begin{equation*}
\varphi_* : \mathrm{CH}(\mathbb{P}(N)) \longrightarrow \mathrm{CH}(\check{\mathbb{P}}^r) \quad 
\end{equation*}

%%%%%%%%%%  scan idx 250  |  folio 243  %%%%%%%%%%

 on a

\begin{equation*}
\varphi_*( \mathbb{P}(N)) = \deg(\varphi)\cdot \check{X} \quad ,
\end{equation*}
de sorte que 5.2.1 se récrit

\begin{equation}\tag{5.2.2}
\deg(\check{X}) = \frac{1}{\deg(\varphi)} \int_{\check{\mathbb{P}}^r} L\cdot\varphi_*( \mathbb{P}(N)) \quad ;
\end{equation}
par abus de notations, dans $\deg(\varphi)$ , $\varphi$ désigne évidemment le morphisme $\mathbb{P}(N) \longrightarrow \check{X}$ induit par (3.1.2), qui est un entier $\geqslant 0$ , nul si et seulement si $\dim X < r-1$ ($=\dim \mathbb{P}(N)$). Par la formule de projection pour $\varphi$ , ceci se récrit

\begin{equation}\tag{5.2.3}
\deg(\check{X}) = \frac{1}{\deg(\varphi)} \int_{\mathbb{P}(N)} \varphi^*(L) \quad .
\end{equation}
Désignons par $\eta \in \mathrm{CH}( \check{\mathbb{P}}^r)$ la classe d'un hyperplan, de sorte que $L = \eta^{r-1}$ , et (5.2.3) se récrit

\begin{equation}\tag{5.2.4}
\deg(\check{X}) = \frac{1}{\deg(\varphi)} \int_{\mathbb{P}(N)} (\varphi^*(\eta))^{r-1} \quad .
\end{equation}
Calculons maintenant la classe $\varphi^*(\eta)$ en termes ``concrets'' .

Notons par $\xi$ (resp. $\tau$) le faisceau inversible canonique $\underline{O}_{\mathbb{P}(N)}(1)$ sur $\mathbb{P}(N)$ (resp. $\underline{O}_{\mathbb{P}(\underline{O}_X(1)^{r+1})}$ sur $\mathbb{P}(\underline{O}_X(1)^{r+1})$ ) . Par fonctorialité, nous avons (cf. (5.1.6))

\begin{equation}\tag{5.2.5}
\alpha^*(\tau) = \xi \quad .
\end{equation}

%%%%%%%%%%  scan idx 251  |  folio 244  %%%%%%%%%%

D'autre part, moyenant l'identification $\mathbb{P}(\underline{O}_X(1)^{r+1}) \simeq X \times \check{\mathbb{P}}^r$ , on a évidemment, posant $H = \underline{O}_X(1)$ ,

\begin{equation}\tag{5.2.6}
\tau = \mathrm{pr}_1^*(H) \otimes \mathrm{pr}_2^*(\eta) \quad .
\end{equation}
Désignons par $\pi : \mathbb{P}(N) \longrightarrow X$ le morphisme structural, on a

\begin{equation*}
\begin{cases} \pi = \mathrm{pr}_1 \circ \alpha \\ \varphi = \mathrm{pr}_2 \circ \alpha \quad , \end{cases}
\end{equation*}
d'où, appliquant $\alpha^*$ aux deux membres de (5.2.6), on trouve, par (5.2.5)

\begin{equation}\tag{5.2.6}
\xi = \alpha^*(\tau) = \pi^*(H) \otimes \varphi^*(\eta) \quad .
\end{equation}
Dans $\mathrm{CH}(\ \mathbb{P}(N))$ , on a donc

\begin{equation}\tag{5.2.7}
\varphi^*(\eta) = \xi - \pi^*(H) \quad ,
\end{equation}
et (5.2.3) se récrit

\begin{equation}\tag{5.2.8}
\deg(\check{X}) = \frac{1}{\deg(\varphi)} \int_{\mathbb{P}(N)} (\xi - \pi^*(H))^{r-1} \quad .
\end{equation}
Le deuxieme membre de 5.2.8 ``se calcule'' grâce au lemme suivant (cf. 5.3.3) :

\underline{Lemme} 5.3. \underline{Soient} $X$ \underline{une variété quasi-projective}, \underline{lisse et connexe de dimension} $n$ , \underline{et} $M$ \underline{un faisceau  localement libre de rang fini} $a$ \underline{sur} $X$ . \underline{Notons par} $\xi$ \underline{le faisceau inversible} $\underline{O}_{\mathbb{P}(M)}(1)$ \underline{sur} $\mathbb{P}(M)$ ,
%%%%%%%%%%  scan idx 252  |  folio 245  %%%%%%%%%%
\underline{et par} $\pi : \mathbb{P}(M) \longrightarrow X$ \underline{le morphisme structural.} \underline{On sait que} $\mathrm{CH}(\ \mathbb{P}(M))$ \underline{est un} $\mathrm{CH}(X)$ \underline{module} (\underline{via} $\pi^*$ ) \underline{qui est libre,} \underline{de base} $1$ , $\xi, \ldots, \xi^{a-1}$ , \underline{où} $\xi$ \underline{satisfait à l'équation}

\begin{equation}\tag{5.3.1}
\xi^a - C_1\xi^{a-1} + C_2\xi^{a-2} - C_3\xi^{a-3} \ldots = 0 \quad ,
\end{equation}
\underline{les} $C_i$ \underline{dénotant les classes de Chern dans} $\mathrm{CH}(X)$ \underline{de} $M$ . \underline{On pose} $C(M) = 1+C_1+C_2+\ldots$, \underline{la classe de Chern totale de} $M$ ; \underline{c'est un élément inversible de} $\mathrm{CH}(X)$ .

\underline{Alors:}

(i) \underline{On a la formule}

\begin{equation}\tag{5.3.2}
\pi_*\!\left(\frac{\xi^{a-1}}{1 + \xi}\right) = \frac{1}{C(M)} \quad .
\end{equation}
(ii) \underline{Soit} $F(T) \in \mathrm{CH}(X)\ [T]$ \underline{un polynôme à coefficients dans} $\mathrm{CH}(X)$ , \underline{tel que le coefficient de} $T^i$ \underline{soit un élément homogène de degré} $n+a-1-i$ . \underline{Donc} $F(\xi)$ \underline{est un élément de degré} \underline{``maximum''} $n+a-1=\dim \mathbb{P}(M)$ \underline{dans} $\mathrm{CH}(\ \mathbb{P}(M))$. \underline{Ceci posé,} \underline{on a}

\begin{equation}\tag{5.3.3}
\int_{\mathbb{P}(M)} F(\xi) = (-1)^{a-1} \int_X \frac{F(-1)}{C(M)} \quad .
\end{equation}
\underline{Démonstration.} Commençons par déduire (5.3.3) de (5.3.2).

Ecrivons

\begin{equation}\tag{5.3.4}
\begin{gathered} \frac{1}{C(M)} = \sum_{j \geq 0} b_j \quad , \quad b_o = 1 \ , \ b_j \ \text{de degré } j \ \text{dans } \mathrm{CH}(X) \\ b_j = 0 \ \text{pour } j < 0 \quad , \end{gathered}
\end{equation}

%%%%%%%%%%  scan idx 253  |  folio 246  %%%%%%%%%%
de sorte que (5.3.2) se récrit

\begin{equation}\tag{5.3.5}
\pi_*(\xi^{a-1+j}) \;=\; (-1)^j b_j \qquad \mbox{pour} \quad a-1+j \geq 0 \ ,
\end{equation}
(compte tenu que $\pi_*(\xi^j) = 0$ si $j < a-1$ pour raison de degrés).

Déduisons-en (5.3.3). Par linéarité, il suffit de traiter le cas $F(T) = uT^{a-1+j}$ , $u$ homogène de degré $n-j$ dans $\mathrm{CH}(X)$ . Par (5.3.5) et la formule de projection pour $\pi$ , on a

\begin{equation}\tag{5.3.6}
\int_{\mathbb{P}(M)} F(\xi) \;=\; \int_{\mathbb{P}(M)} \pi^*(u) \cdot \xi^{a-1+j} \;=\; \int_X u \cdot \pi_*(\xi^{a-1+j}) \;=\; (-1)^j \int_X u \cdot b_j \ .
\end{equation}
D'autre part,

\begin{equation}\tag{5.3.7.}
(-1)^{a-1} \int_X \frac{F(-1)}{C(M)} \;=\; (-1)^{a-1} \int_X \frac{(-1)^{a-1+j}u}{C(M)} \;=\; (-1)^j \int_X \sum_i \, u \, b_i \ ,
\end{equation}
et compte tenu de ce que $\deg(u) = n-j$ , et $\deg(b_i) = i$ , on a

\begin{equation*}
\int_X u \cdot b_i = 0 \qquad \mbox{si} \quad i \neq j \ ,
\end{equation*}
de sorte que (5.3.7) se récrit

\begin{equation}\tag{5.3.8}
(-1)^{a-1} \int_X \frac{F(-1)}{C(M)} \;=\; (-1)^j \int_X u \cdot b_j \ ,
\end{equation}
d'où l'implication annoncée.

Démontrons maintenant (5.3.2), ou se qui revient au même, (5.3.5).


%%%%%%%%%%  scan idx 254  |  folio 247  %%%%%%%%%%

On sait que

\begin{equation*}
\pi_*(\xi^j) = 0 \qquad \mbox{si} \quad j < a-1
\end{equation*}
pour une raison de degré, et que

\begin{equation*}
\pi_*(\xi^{a-1}) \;=\; 1_X
\end{equation*}
pour une raison géométrique ( $\xi^{a-1}$ rencontre chaque fibre "en un point"), de sorte que, par la formule de projection, on a

\begin{equation}\tag{5.3.9}
\pi_*\left( \sum_{j=0}^{a-1} \pi^*(U_i) \, \xi^i \right) \;=\; U_{j-i} \ .
\end{equation}
Par (5.3.9), l'énoncé 5.3.5 est conséquence du lemme suivant  dans lequelle $\xi$ devient $T$ , et $C_i$ devient $S_i$, compte tenu de l'équation 5.3.1.

\underline{Lemme} 5.4. \underline{Soit} $a$ \underline{un entier $> 0$} , \underline{et désignons par} $A$ \underline{l'anneau} $Z[[X_1,\ldots,X_a]]$ \underline{des séries formelles en} $a$ \underline{indéterminées.} \underline{Désignons par} $B$ \underline{la} $A$-\underline{algèbre}

\begin{equation*}
B = A[[T]] \, / \, \left( \prod_{i=1}^{a} (T-X_i) \right) \ ,
\end{equation*}
\underline{de sorte que} $B$ \underline{est un} $A$-\underline{module libre} , \underline{à base} $1,T,\ldots,T^{a-1}$ ,

\begin{equation}\tag{5.4.1}
B = A \, T^{a-1} \oplus A \, T^{a-2} \oplus \ldots \oplus A \ .
\end{equation}

%%%%%%%%%%  scan idx 255  |  folio 248  %%%%%%%%%%

\underline{Désignons par} $S_i \in A$ \underline{la $i$-ième fonction symétrique}

\begin{equation*}
S_i = \sum_{1 \leq \alpha_1 < \ldots < \alpha_i \leq a} X_{\alpha_1} \ldots X_{\alpha_i} \ ,
\end{equation*}
\underline{de sorte que}

\begin{equation}\tag{5.4.2}
\sum_{i=0}^{a} (-1)^i \, S_i \, T^{a-i} = 0 \quad \mbox{dans } B \ .
\end{equation}
\underline{L'élément} $1 + S_1 + S_2 + \ldots + S_a$ \underline{est inversible dans} $A$ , \underline{nous écrivons son inverse}

\begin{equation}\tag{5.4.3}
\frac{1}{\sum\limits_{i=0}^{a} S_i} = \sum_{j \geq 0} b_j \ , \qquad b_j \ \mbox{homogène de degré } j \mbox{ dans } A \ , \quad b_o = 1 \ .
\end{equation}
\underline{D'autre part, pour tout entier} $j \geq 0$ \underline{on définit un élément} $a_j \in A$ \underline{homogène de degré} $j$ \underline{en écrivant, par} 5.4.1

\begin{equation}\tag{5.4.4}
T^{a-1+j} = a_j \, T^{a-1} + ? \, T^{a-2} + ? \, T^{a-3} + \ldots + ? \ , \quad \mbox{les } ? \in A \ .
\end{equation}
\underline{Ceci posé, on a pour tout} $j \geq 0$

\begin{equation}\tag{5.4.5}
a_j = (-1)^j \, b_j \ .
\end{equation}
\underline{Démonstration}. Compte tenu de la définition (5.4.3) des $b_j$ , tout revient à démontrer qu'on a 


%%%%%%%%%%  scan idx 256  |  folio 249  %%%%%%%%%%

\begin{equation}\tag{5.4.6}
 \begin{array}{l} a_o = 1 \\[4pt] \sum\limits_{i=0}^{a} (-1)^i \, S_i \, a_{j-i} = 0 \qquad \mbox{si } j > 0 \ . \end{array}
\end{equation}
Or, on a évidemment $a_o = 1$ , et la deuxième relation (pour $j > 0$ donné) n'est autre que celle déduite ,en prenant le coefficient de $T^{a-1}$ (via 5.4.1), de la relation

\begin{equation}\tag{5.4.7}
\sum_{i=0} (-1)^i \, S_i \, T^{a-1+j-i} = 0 \quad \mbox{dans } B \ ,
\end{equation}
relation obtenue en multipliant 5.4.2 par $T^{j-1}$ , cqfd.

\underline{Corollaire} 5.5. \underline{On a la formule}

\begin{equation}\tag{5.5.1}
\deg(\check{X}) = \frac{(-1)^n}{\deg(\varphi)} \int_X \frac{(1+H)^{r-1}}{C(N)} \ .
\end{equation}
\underline{Démonstration}. On applique (5.3.3) (pour $M = N$ , le faisceau normal de $X \hookrightarrow \mathbb{P}^r$) à (5.2.8).

Ecrivons, pour une variété quasi-projective et lisse $S$ , sa classe de Chern totale (cf. 5.3)

\begin{equation}\tag{5.5.2}
C(S) \stackrel{\mathrm{dfn}}{=} C(T_S)
\end{equation}
où $T_S$ désigne le fibré \underline{tangent} de $S$ .


%%%%%%%%%%  scan idx 257  |  folio 250  %%%%%%%%%%

La duale de la suite exacte (5.1.2) montre qu'on a

\begin{equation}\tag{5.5.3}
C(\mathbb{P}^r) \;=\; (1 + H)^{r+1} \qquad \text{dans } CH(\mathbb{P}^r) \quad .
\end{equation}
La duale de la suite exacte(5.1.1)montre qu'on a

\begin{equation}\tag{5.5.4}
C(N)\, C(X) = C(\mathbb{P}^r)|X \qquad \text{dans } CH(X) \quad .
\end{equation}
Mettant ensemble (5.5.3) et (5.5.4),(5.5.1)se récrit

\underline{Corollaire} 5.6. \qquad \underline{On a la formule}

\begin{equation}\tag{5.6.1}
\deg(\check{X}) \;=\; \frac{(-1)^r}{\deg(\varphi)} \; \int_X \frac{C(X)}{(1+H)^2} \quad .
\end{equation}
(5.7) Pour une variété $S$ projective et lisse, on définit sa \underline{caractéristique d'Euler-Poincaré par}

\begin{equation}\tag{5.7.1}
\chi(S) \;=\; \int_S C(S) \quad .
\end{equation}
\underline{Proposition} 5.7.2. \qquad \underline{Désignons par} $Y$ \underline{une section hyperplanelisse de} $X$ , \underline{et par} $\Delta$ \underline{une section lisse de} $X$ \underline{par un sous-espace linéaire de} $\mathbb{P}^r$ \underline{de codimension deux.} \underline{On a}

\begin{equation}\tag{5.7.3}
\deg(\check{X}) = \frac{(-1)^n}{\deg(\varphi)} \quad \left[ \chi(X) + \chi(\Delta) - 2\chi(Y) \right] \quad .
\end{equation}

%%%%%%%%%%  scan idx 258  |  folio 251  %%%%%%%%%%

\underline{Démonstration.} \quad Notons $i : Y \hookrightarrow X$ l'inclusion. On a

\begin{equation*}
N_{Y \longrightarrow X} \;=\; H|Y \;=\; i^{*}H \quad ,
\end{equation*}
de sorte que la suite exacte

\begin{equation*}
0 \longrightarrow T_Y \longrightarrow i^{*}T_X \longrightarrow i^{*}H \longrightarrow 0
\end{equation*}
donne

\begin{equation}\tag{5.7.4}
C(Y) \;=\; i^{*} \left( \frac{C(X)}{1+H} \right) \quad .
\end{equation}
On a donc

\begin{equation*}
\chi(Y) \;=\; \int_Y i^{*} \left( \frac{C(X)}{1+H} \right)
\end{equation*}
et, en appliquant la formule de projection

\begin{equation}\tag{5.7.5}
\chi(Y) \;=\; \int_X i_{*}i^{*} \left( \frac{C(X)}{1+H} \right) \;=\; \int_X \frac{H \cdot C(X)}{1 + H} \quad .
\end{equation}
Itérant ce calcul, on trouve pour $\Delta$

\begin{equation}\tag{5.7.6}
\chi(\Delta) \;=\; \int_Y \frac{H\ C(Y)}{1 + H} \;=\; \int_Y i^{*} \left( \frac{H}{(1+H)} \cdot \frac{C(X)}{(1+H)} \right) \;=\; \int_X \frac{H^2\ C(X)}{(1+H)^2} \quad ,
\end{equation}
de sorte que

\begin{equation*}
\begin{aligned} \chi(X) - 2\chi(Y)+\chi(\Delta) &= \int_X C(X) \left[ 1 - \frac{2\ H}{1+H} + \frac{H^2}{(1+H)^2} \right] \\ &= \int_X \frac{C(X)}{(1 + H)^2} \quad , \end{aligned}
\end{equation*}
ce qui achève la démonstration, compte tenu de 5.6.1.


%%%%%%%%%%  scan idx 259  |  folio 252  %%%%%%%%%%


6. \underline{Le cas d'une base générale}

6.1. Soient $S$ un schéma et $\bar{s}$ un point géométrique de $S$, défini par une clôture algébrique du corps résiduel $k(s)$ pour $s \in S$. Désignons par $\mathbb{P}^r$ l'espace projectif à $r$ dimensions sur $S$ , par $\check{\mathbb{P}}^r$ l'espace projectif dual et par $I \subset \mathbb{P}^r \times_S \check{\mathbb{P}}^r$ la variété d'incidence. Soit $D \subset \check{\mathbb{P}}^r$ une droite, définie par un point à valeurs dans $S$ d'une grassmannienne convenable, et $\Delta$ l'axe du pinceau d'hyperplans correspondant (cf. 2.1).

Soit $X$ un $S$-schéma propre et lisse, purement de dimension $n \geq 1$, muni d'un plongement projectif $X \hookrightarrow \mathbb{P}^r$ . Soient $\widetilde{X} = (X \times_S D) \cap I$ le schéma sur $D$ de fibres les sections hyperplanes de $X$ paramètrées par $D$, et $u$ la projection de $\widetilde{X}$ sur $D$.

Notons par un indice $\bar{s}$ le changement de base par $\bar{s} \longrightarrow S$.

Supposons que $\Delta_{\bar{s}}$ soit transverse à $X_{\bar{s}}$ . Alors, au-dessus d'un voisinage $U$ de $s$ dans $S$, $\Delta$ est transverse à $X$ . On suppose dans ce qui suit que $U = S$ . Le schéma $\widetilde{X}$ est alors lisse sur $S$ , et un voisinage de l'image réciproque de $\Delta$ dans $\widetilde{X}$ est lisse sur $D$ .

Supposons de plus qu'une des sections hyperplanes de $X_{\bar{s}}$ appartenant au pinceau $D_{\bar{s}}$ soit lisse. Alors, après que l'on ait remplacé $S$ par un voisinage convenable de $s$, le sous-schéma $F$ de $\widetilde{X}$ défini par l'idéal jacobien $J^{n-1}(\widetilde{X}/D)$ est fini et plat sur $S$ .


%%%%%%%%%%  scan idx 260  |  folio 253  %%%%%%%%%%

6.2. \underline{Le lieu exceptionnel} du pinceau $D$ est le diviseur de Cartier relatif $E$ de $D/S$ défini par le $\mathcal{O}_D$-Module $u_*\mathcal{O}_F$, fini et plat sur $S$ . Localement sur $D$, si $\mathcal{O}_D^n \xrightarrow{\ v\ } \mathcal{O}_D^n \longrightarrow u_*\mathcal{O}_F \longrightarrow 0$ est une présentation de $u_*\mathcal{O}_F$ , $E$ a pour équation $\det(v) = 0$. La formation de $E$ est compatible à tout changement de base.

Le diviseur $E_{\bar{s}}$ de $D_{\bar{s}}$ s'écrit $\Sigma\, n(P).P$, pour $n(P)$ la somme des colongeurs de l'idéal jacobien $J^{n-1}(\widetilde{X}/D)$ en les points critiques de $u$ d'image $P$ .

\underline{Proposition} 6.3. \underline{Si} $D_{\bar{s}}$ \underline{est un pinceau de Lefschetz de sections hyperplanes de} $X_{\bar{s}}$ , \underline{et si} car $k(s) \neq 2$ \underline{ou que} $n$ \underline{est pair}, \underline{il existe un voisinage} $U$ \underline{de} $s$ \underline{dans} $S$ \underline{tel que}

(a) \underline{Pour tout point géométrique} $\bar{t}$ \underline{de} $U$, $D_{\bar{t}}$ \underline{est un pinceau de Lefschetz de sections hyperplanes de} $X_{\bar{t}}$

(b) \underline{Au-dessus de} $U$, \underline{le lieu exceptionnel} $E$ \underline{est étale sur} $S$.

Pour prouver (a) et (b), il suffit de noter que, sous les hypothèses faites, $D_{\bar{s}}$ vérifie la condition c) de 2.2 si et seulement si le diviseur $E_{\bar{s}}$ est réduit. On pourrait, en invoquant 3.3, n'utiliser que la moitié la plus facile $(\Longrightarrow)$ de cette équivalence.

6.3. La proposition devient fausse dans le cas d'exception cité. Toutefois, si $S$ est \underline{de caractéristique} 2 et que $n$ est impair, il existe un voisinage $U$ , de $S$ au-dessus duquel (a) est vrai, et on peut définir un diviseur $E'$ étale sur $U$ tel que $E = 2 E'$. Nous ne ferons pas usage de ce résultat.


\underline{Références}

[1] GROTHENDIECK, A. \quad La théorie des Classes de Chern, \\ \hspace*{5.5em} Bull. Soc. Math. France, 86, 1958, 137-154.
